\section{The determinant-to-entry map and the finite exceptional corner}
\label{sec:det-entry-finite-corner}

We now identify the determinant ring and the entry
space used in the two propagation runs.
This also constructs the finite weight-space ideal
covering every paired exceptional component.
All rings and representations in this section are
the actual objects constructed in
Sections~\ref{sec:entry-space-flat} and
\ref{sec:entry-nice-sharp}.

\subsection{Traces generate the weight-zero coefficient algebra}
Let \(S_n=(B_n,A_n,\rho_n,\psi_n)\), \(n=2,3\).
Suppose a complete local coefficient algebra \(C\)
maps to the ambient diagonal deformation ring,
and its image contains every
\(\operatorname{tr}\rho_n(g)\) and every
\(\psi_{vi}(s)\).
We claim that its image contains \(B_n\).

The image is a complete local quotient of \(C\)
with the same residue field.
It is closed in the ambient ring.
Indeed, at each fixed precision the induced and
intrinsic maximal-ideal topologies are cofinal
on a finite entry lattice, as in (LS3);
a compatible approximation therefore has a limit
in the image, and separatedness identifies its
image with the given limit.

Write
\[
 X=\rho_n(g_0)-\rho_n(g_0^*)
                  =\zeta\operatorname{diag}(w_i).
\]
Here \(\zeta\) is a unit with a prescribed
nonzero residue, and
\[
                  \operatorname{tr}(X^2)=2\zeta^2.            \tag{DC1}
\]
The trace of every product of group matrices
belongs to the coefficient image.
Thus \(\zeta^2\) belongs to it.
Hensel lifting \(Z^2-\zeta^2\), using the
specified residue of \(\zeta\) and the unit
\(2\zeta\), shows that \(\zeta\) itself belongs
to the image.
Consequently the coordinate projectors belong
to the coefficient span of group matrices:
\[
 E_+=\tfrac12(W^2+W),\quad
 E_-=\tfrac12(W^2-W),\quad
 E_0=1-W^2,\qquad W=\zeta^{-1}X,                              \tag{DC2}
\]
where \(E_0\) is omitted in rank two.
The span is a subalgebra.
It follows that
\[
 \rho(g_1)_{i_1i_2}\cdots
 \rho(g_r)_{i_ri_1}
   =\operatorname{tr}
       (E_{i_1}\rho(g_1)E_{i_2}\rho(g_2)
                    \cdots E_{i_r}\rho(g_r))                 \tag{DC3}
\]
belongs to the same coefficient image.

For completeness this captures every
weight-zero monomial, including ones whose
indices initially do not form a cycle.
In rank two the only nonzero entry weights
are \(2,-2\).
In rank three they are \(1,2,-1,-2\);
polarization identifies the two weight-one
entry spans, and also the two weight-minus-one
entry spans.
Every primitive multiset of these nonzero
weights with sum zero is one of
\[
                 (1,-1),\quad(2,-2),\quad
                        (2,-1,-1),\quad(-2,1,1).
\]
After the indicated reflection of an adjacent
entry, each is a closed index cycle.
Equation (DC3) applies.
Remove such a primitive submultiset and
induct; diagonal factors already belong to
the image by (DC2).
Finally take the defining adic closure of
the weight-zero algebra.
This proves the claim, including at \(\ell=3\):
only \(2\) and \(\zeta\) were inverted.

Let \(\mathcal B_n\) be the actual universal
ordinary determinant ring for the corresponding
scalar-twisted arithmetic packet.
The entry representation, composed with the
surjection of arithmetic groups and twisted
by the fixed scalar \(\kappa\), gives its
continuous determinant and local characters.
It satisfies its residual, polarization, and
ordered-product relations.
At an inertia generator killed in the
arithmetic quotient, a prime-to-\(\ell\)
power is one and its reduction is one.
The geometric sum in
\[
                  M^m-1=(M-1)(1+M+\cdots+M^{m-1})
\]
is invertible, since its reduction is
\(mI\); hence \(M=1\).
Thus these are the exact unramified
relations as well.

Universality gives a map
\[
                       \Phi_n:\mathcal B_n\longrightarrow B_n.
                                                               \tag{DC4}
\]
The image contains the entry traces:
the arithmetic group map is surjective,
and the scalar \(\kappa(g)\) is a unit
of \(\mathcal O\).
It contains the local characters by the
surjectivity of the local group maps.
The preceding claim makes (DC4) surjective.
This proves the source's \srcid{PhiEx},
including the asserted target ring.

\subsection{The actual rank-two corner on the exceptional quotient}
Write \(\mathfrak e\subset A_3\) for the
odd-entry ideal, and put
\[
 C=B_3/(\mathfrak e\cap B_3),\qquad
                 Q=A_3/\mathfrak e .
\]
If \(\mathfrak e=A_3\), take \(I_3=B_3\)
below; there are no exceptional primes.
Otherwise \(C\hookrightarrow Q\), and
\(C\) is complete local Noetherian.
The universal representation over \(Q\) has
block form
\[
 \rho(g)=
 \begin{pmatrix}
 a(g)&0&b(g)\\
 0&\chi(g)&0\\
 c(g)&0&d(g)
 \end{pmatrix},\qquad
       \sigma(g)=\begin{pmatrix}a(g)&b(g)\\c(g)&d(g)\end{pmatrix}.
                                                               \tag{DC5}
\]
The outer block is a representation.
Its pairing is the rank-two reversal pairing,
and its marked normalization is unchanged.
Its local characters are
\(\psi_{v0},\psi_{v2}\).

We check its shorter ordinary identity rather
than canceling a matrix factor without cause.
At each \(v\) choose the frame's element \(u_v\)
with \(\Psi(u_v)^2\ne1\).
The three quantities
\[
 \chi(u_v)-\psi_{v0}(u_v),\quad
 \chi(u_v)-\psi_{v2}(u_v),\quad
                   \det(\sigma(u_v)-\psi_{v1}(u_v)I_2)
\]
are units.
Their residues are respectively
\(1-\Psi(u_v)\), \(1-\Psi(u_v)^{-1}\),
and
\((\Psi(u_v)-1)(\Psi(u_v)^{-1}-1)\).
The corner product \(b(u_v)c(u_v)\) has
weight zero and lies in the maximal ideal
of \(C\), so it does not alter the last
residue.

For \(P=\sigma(x)-\psi_{v0}(x)I_2\),
\(M=\sigma(u_v)-\psi_{v1}(u_v)I_2\),
and \(R_y=\sigma(y)-\psi_{v2}(y)I_2\),
the rank-three identity gives \(PMR_y=0\)
for every \(y\).
Since \(\psi_{v2}\) is a character,
\[
 PM\sigma(y)=\psi_{v2}(y)PM,\qquad
 PM^2=(\psi_{v2}(u_v)-\psi_{v1}(u_v))PM.
\]
The two-by-two Cayley--Hamilton identity gives
\[
 (\det M)P=
  \bigl(\operatorname{tr}M
        -\psi_{v2}(u_v)+\psi_{v1}(u_v)\bigr)PM.
\]
Multiply by \(R_y\) and use the unit
\(\det M\): \(PR_y=0\).
Similarly the middle entry of the original
product at \((u_v,g,u_v)\) gives
\(\chi(\operatorname{loc}_v g)=\psi_{v1}(g)\).
Thus (DC5) is the actual normalized
ordinary rank-two corner.

The entry spans of \(\sigma\) are finite
\(C\)-modules, their adic continuity descends
from the rank-three spans, and every balanced
off-diagonal monomial belongs to
\(\mathfrak m_C\).
The vertex construction (LS3) applies with
coefficient ring \(C\), target \(Q\), and these
finite spans: the graded Laurent subalgebra
is Noetherian of finite type over \(C\), its
vertex completion has diagonal residue, and
evaluation at its Laurent variable \(1\)
recovers the target matrices.
Its degree-zero image lies in the closed
copy of \(C\).
We therefore obtain actual maps
\[
           \xi:A_2\longrightarrow Q,\qquad
             \beta:B_2\longrightarrow C,\qquad
      \xi|_{B_2}=\iota\beta,                                 \tag{DC6}
\]
with the outer entries and characters of
(DC5).

\subsection{Construction of the finite exceptional ideal}
Let \(I_2\subset B_2\) be the supported ideal
after the rank-two propagation run.
The conclusions already proved for that run
are
\[
 \mathfrak m_{B_2}\subset
               \sqrt{I_2+\mathfrak l_2},
 \quad
 I_2\subset P_2\cap B_2
           \ \text{at every rank-two Type-I prime }P_2.       \tag{DC7}
\]
Here \(\mathfrak l_n\) is the source weight
augmentation: it contains
\(\mathfrak m_{\mathcal O}\) and
\(\psi_{vi}(g)-1\) for
\(g\in I_v,\ \Psi(\operatorname{loc}_v g)=1\).
Define, with \(p:B_3\twoheadrightarrow C\),
\[
                         I_3=p^{-1}(\beta(I_2)C).             \tag{DC8}
\]
This is a specified ideal in the same
rank-three coefficient ring.

At a Type-E prime \(P_3\subset A_3\), the
outer block stays off both triangular
directions, since its two corner entries
are the only possible off-diagonal entries
left in (DC5).
Thus \(\xi^{-1}(P_3/\mathfrak e)\) is
rank-two Type I.
Equation (DC7), followed through (DC6),
gives
\[
                           I_3\subset P_3\cap B_3.            \tag{DC9}
\]

We prove its finiteness over the weight
augmentation.
Let \(P\subset B_3\) be any prime containing
\(I_3+\mathfrak l_3\).
It contains \(\ker p\), and (DC6) pulls
\(\mathfrak l_3\) back to
\(\mathfrak l_2\).
The prime of \(B_2\) so obtained contains
\(I_2+\mathfrak l_2\); hence (DC7) makes
every element of \(B_2\), modulo this
prime, a coefficient-field constant.

Modulo \(P\), the middle function \(\chi\)
is a character, is \(1\) residually, and
satisfies \(\chi(g^*)=\chi(g)\).
Its inertia restrictions on
\(I_v\cap\ker\Psi\) are one by the weight
augmentation and the equality
\(\chi|_{G_v}=\psi_{v1}\).
Finite prime-to-\(\ell\) inertia quotients
contribute nothing to a principal-unit
character in characteristic \(\ell\).
The arithmetic abelian control proved
with the frame therefore gives a fixed
nonzero exponent \(h\) with
\(\chi(g)^h=1\) for all \(g\).
This is exactly its finite class-group
exponent, and is independent of \(g\).

A principal unit in a characteristic-\(\ell\)
local domain has no nontrivial finite
order: write \(h=\ell^a m\) with
\(\ell\nmid m\); the \(\ell^a\)-power
reduces to \(x^m=1\), and
\(1+x+\cdots+x^{m-1}\) is a unit.
Hence
\[
                              \chi(g)=1\pmod P.              \tag{DC10}
\]
Every weight-zero rank-three entry monomial
modulo \(\mathfrak e\) is either zero or
a product of middle entries and a
weight-zero outer-block monomial.
Equations (DC6), (DC7), and (DC10) make
it a coefficient-field constant modulo
\(P\).
The outer local characters are constant
by (DC6), and the middle local character
equals the middle global character.
Thus all the topological generators of
\(B_3\) are coefficient-field constants
modulo \(P\).
The induced topology on its degree-zero
part is its maximal-ideal topology by
(LS2); the ideal \(P\) is closed.
It follows that \(B_3/P\) is the
coefficient residue field, or equivalently
\(P=\mathfrak m_{B_3}\).
We have proved
\[
 \mathfrak m_{B_3}\subset
                     \sqrt{I_3+\mathfrak l_3}.               \tag{DC11}
\]

Equations (DC9)--(DC11) construct the
finite exceptional ideal required by
the rank-three component argument.
In particular it is obtained from
the actual rank-two supported quotient;
it is not an independent finiteness
assumption on all exceptional points.

\paragraph{Exact source correspondence.}
The trace calculation is
\srcid{FrobeniusSmoothVectors.TraceMem.B\_gen};
the group quotient, scalar twist, continuity,
and surjection (DC4) are
\srcid{CurveDVRSpecialization.PxIsCont.phiEx}.
The corner calculation expands
\srcid{ExistsFrobeniusCorner},
\srcid{FrobeniusCornerOrderSkips},
\srcid{FrobeniusCornerUnits}, and
\srcid{FrobeniusCornerMemCone}.
The middle-character argument is
\srcid{FrobeniusCornerMiddleTrivial} and
\srcid{FrobeniusCornerMiddleMem}.
The finite ideal (DC8) and its two
properties are
\srcid{FrobeniusSmoothVectors.%
QuotientHomNonvanishingBound.l53}.
These are frozen-source modules under
\srcid{Development.AutomorphicArithmetic.AutomorphicVarieties}
at tag \srcid{alpoge-source-7cbf7a0c}.
